\scanpage{76}
% archive sheet number: 75

part, \inserted{\handwritten{\uline{une fois choisi} \uline{une clôture} $\overline{\mathbb{F}}_p$,}} la gerbe $\handwritten{\unsure{\underline{\pi}|\underline{\Pi}}}_p(\underline{F}_p)$, grâce àu\deleted{son} foncteur \unsure{fvbre|fibre|fïbre} évident ``fibre en $\handwritten{\bar{\underline{F}}}_p$ '' sur la catégorie des $\underline{Q}_p$-faisceaux constructibles sur $\underline{F}_p$, s'identifie à la \deleted{catégorie}\inserted{gerbe triviale définie\deleted{\unsure{et}} par l'}enveloppe $\underline{Q}_p$-\unsure{algébriqueque|algébrique} du groupe profini $\handwritten{\pi}_1(\underline{F}_p)=\handwritten{\bar{\underline{Z}}}\deleted{\unsure{r}}$, \deleted{lequel, comme on le constate aisément, est le}

Un \inserted{$\underline{Q}_p$-}faisceau constructible sur $\underline{F}_p$ s'identifie à un vectoriel de dimension finie $E$ sur $\underline{Q}_p$, muni d'un automorphisme $F_E$ tel que l'\deleted{hom} opération $n\handwritten{\mapsto}F_E^n$ de $\underline{Z}$ sur $E$ se prolonge par continuité en une opération de \deleted{\unsure{$\hat{\underline{Z}}$}} $\handwritten{\hat{\underline{Z}}}$ \margin{\handwritten{$Z$}}. Identifiant \inserted{\handwritten{(par un foncteur $j$, \uline{distinct} de celui de 5.2)}} la catégorie de ces $(E,\deleted{\unsure{E_E}}F_E)$ à une\handwritten{\,}sous-catégorie pleine de\deleted{\unsure{s}} la catégorie \deleted{des} $\mathrm{Fcriso}(\underline{Q}_p)$ des couples $(E,F_E)$, où $F_E$ est un automorphisme \uline{quelconque}, on voit que le lien \deleted{de notre}\inserted{$G\handwritten{'}$ de} $\handwritten{\underline{\pi}}_p(\underline{F}_p)$ de cette catégorie s'identifie à un quotient du lien $G$ de $\mathrm{Fcriso}(\underline{F}_p)$ déterminé plus haut, et on voit aussitôt que ce quotient est de la forme $G_u\mathrm{x}G_r^{\deleted{\unsure{x}}\handwritten{'}}$, où $G_r'$ est \unsure{le|he} qutient de $G_r$ (ceci revient à dire que si on décompose $F_E$ en produit de sa composante unipotente par sa composante semi-simple, alors la condition que $F_E$ définisse une opération de $\handwritten{\hat{\underline{Z}}}$ ne dépend que de la composante semi-simple); de façon précise, $G'_r$ est le groupe quotient de $G_r$ défini par le sous-groupe du groupe des caractères $\handwritten{\bar{\underline{Q}}}_p^{\handwritten{*}}$ de $G_r$ formé des \uline{unités} p\uline{\unsure{éadiques|-adiques}} : cela signifie en effet que $F_E$ définit uneopération de \deleted{\unsure{$F_E$}} $\handwritten{\hat{\underline{Z}}}$ (i.e. définit un F-isocristal de pente nulle) sss les valeurs\deleted{\unsure{s}} propres de $F_E$ sont des unités, ce qui résulte en effet aussitôt du critère donné plus haut des F-isocristaux de pente nulle (ou se vérifie directement de façon immédiate; la référence au résultat sur les fourbis de pente nulle devrait se justifier par un grain de sel, car l'inclusion utilisée ici des $\underline{Q}_p$-faisceaux constructibles sur $\underline{F}_p$ dans\deleted{\unsure{à}} les F-iso\deleted{\unsure{cr}}cristaux n'est pas celle utilisée \deleted{plus haut}\inserted{\handwritten{dans 5.2}}, cf. plus bas). On trouve ainsi une structure d'extension sur $G\handwritten{=}\operatorname{Lien}(\underline{G}(\underline{F}_p))$:
\[
0\handwritten{\longrightarrow}H\handwritten{\longrightarrow}G\handwritten{\longrightarrow}G'\handwritten{\longrightarrow}0\ ,
\]
où $H$ est le groupe de t.m. dont le groupe des caractères est le
